\section{Conclusiones}
\label{sec:conclusiones}

Se determinó que la amenaza cuántica no es únicamente un problema futuro: debido al riesgo \textit{store now, decrypt later}, la información cifrada hoy que deba mantenerse confidencial durante años ya se encuentra expuesta. Asimismo, se corroboró que disponer de estándares no equivale a estar preparado: aunque ML-KEM, ML-DSA y SLH-DSA ya están definidos, la dificultad principal de la transición es de ingeniería y gestión, pues exige inventariar dependencias, mantener la interoperabilidad y evitar ataques de \textit{downgrade}.

Se evidenció también que la adopción es desigual: el acuerdo de claves poscuántico ya protege más de la mitad del tráfico humano de Cloudflare, mientras que solo alrededor del 3.7\,\% de los servidores de origen lo soporta \cite{westerbaan2025postquantum}. La demostración realizada permitió comprobar una de las causas de este rezago: una firma ML-DSA-65 ocupa 3309 bytes, unas 50 veces más que una firma Ed25519, lo que encarece la migración de certificados.

A partir del impacto observado en 2026, se concluye que el tema se dirige de la confidencialidad hacia la autenticación y hacia infraestructura más allá de la Web, como IPsec \cite{goldberg2026ipsec,valenta2026authentication}; el incidente de junio de 2026 muestra que la próxima dificultad será preservar la disponibilidad. Se identificó que a la academia le falta una definición y una métrica comunes de cripto-agilidad, y a la industria, criterios verificables para determinar cuándo una organización está preparada. Por ello, un profesional de seguridad debería vigilar la llegada de certificados poscuánticos a las PKI públicas, las hojas de ruta PQC de sus proveedores, la configuración de la negociación híbrida y el inventario criptográfico de su propia organización. Como limitación, la evidencia de despliegue proviene de un solo proveedor, por lo que no es generalizable.
